% =====================================================================
\appendix

\chapter{补充证明与论证}\label{app:proofs}

本附录给出正文中若干结果的完整证明或更细致的论证。凡正文中已给出完整证明者，
此处不重复。

\section{命题 \ref{prop:overconstrain} 的说明}

正文给出的证明基于横截性定理。这里补充两点。

\paragraph{为什么"独立约束"是必要的。}
若两个约束不独立（例如第二个是第一个的推论），则它们对余维数的贡献不叠加，
$\mathrm{codim}\,\mathcal{S} < c$，解集维数高于 $d - c$。因此命题
\ref{prop:overconstrain} 的结论应在"约束独立且横截"的前提下理解。这一前提在底层
设计中通常成立（每条约束来自不同的既有对象：二进制兼容性、调试器、异常展开、
系统调用接口），而在高层设计中常常不成立（多个"约束"实际上是同一偏好的不同表述）。

\paragraph{为什么"解集退化"不等于"唯一解"。}
$c \ge d$ 时解集是有限集合或空集，有限集合可以有多个元素。底层之所以在实践中常常
"接近唯一"，是因为除了形式约束之外还有隐式的经济约束：同时维护两套不兼容的调用
约定会带来人力与工具链成本，这一成本在可选方案之间形成严格的偏好序。

\section{定理 \ref{th:integration} 的完整论证}

正文给出的是主论证。这里补充成本函数的选择与假设的合理性。

\paragraph{成本的定义。}
设单位交互成本为 $c$，指"正确处理一次特定实体与特定资源之间的交互"所需的全部
工作（设计、实现、测试与长期维护）。在无基座情形下，每个实体必须独立完成这一工作，
且必须独立处理资源的不变量。

\paragraph{一致性义务的量级。}
对资源 $r$，$n$ 份独立实现必须在 $r$ 的不变量上取得一致。若一致性不被强制，
则任一份实现的偏离都会破坏全局不变量。若被强制，则强制本身需要 $n$ 个位置都知道
该不变量，即强制点数为 $\Omega(n)$。这两种情形都不优于"读集 $\ge k$"的结论
（定理 \ref{th:centralise}）。换言之，无基座不仅成本更高，而且它把不变量强行散落，
从而\emph{永久提高}了系统的读集下界。这是统一基座的第二个收益，独立于摊销。

\paragraph{假设 (b) 的松弛。}
假设 (b)（每个实体参与全部 $m$ 个资源）是最坏情形。若实体 $e$ 仅参与
$|\mathcal{R}(e)|$ 个资源，则无基座成本为 $\sum_e |\mathcal{R}(e)|$。但只要
$\mathbb{E}_e[|\mathcal{R}(e)|]$ 随 $m$ 增长（即资源的使用是分散而广泛的），
乘积形式的增长仍然成立。因此式 \eqref{eq:threshold} 的定性结论不依赖最坏情形假设。

\section{定理 \ref{th:six} 的逐项论证补充}

正文以"缺失导致未定义行为"的方式逐项论证。这里补充两点对论证强度的说明。

\paragraph{(F3) 调度的论证强度。}
该论证依赖"$\to$ 应当是函数，或其非确定性应被显式规定"。若一个系统的语义被定义为
\emph{所有}交错的集合（即把非确定性当作语义的一部分而不加选择），那么"选择函数"
确实可以缺失。但这种情形下，系统的\emph{实现}仍然必须在每个具体执行中选择一个交错，
因此调度职能在实现层面仍然存在，只是没有出现在抽象语义中。本文的表述"或者在该
系统自身的代码中，或者在其宿主中"覆盖了这一情形。

\paragraph{(F6) 隔离的论证强度。}
该论证依赖"故障可达"。若一个系统的运行环境保证无故障（例如纯函数式计算、
或在被验证为无故障的受限环境中执行），则 (F6) 可以缺失或退化为平凡的。
这正是一些极简系统能够不拥有完整六职能的原因，也说明定理 \ref{th:six} 的前提不是
多余的。

\section{定理 \ref{th:centralise} 的读集下界}

正文证明了 $|\rho(c)| \ge |\mathrm{S}(\iota)|$。这里说明为什么这一不等式是紧的。

若 $\mathrm{S}(\iota) = \{M_0\}$，且 $M_0$ 的接口是充分的（定义
\ref{def:neardecomp}），则任何可能破坏 $\iota$ 的变更都必须在 $M_0$ 的接口上被
表达，因此变更者只需理解 $M_0$ 与其调用点。此时
$|\rho(c)| = 1 + O(\text{调用点数})$，而下界 $1$ 被达到（当变更本身就发生在 $M_0$
内部时）。

若 $\mathrm{S}(\iota) = \{M_0, M_1\}$ 且这两个模块都对 $\iota$ 有部分责任，
则变更者必须理解两者\emph{以及它们之间的一致性规则}，后者的理解成本严格高于前两者
之和。这说明读集关于 $|\mathrm{S}(\iota)|$ 的增长实际上是超线性的，而不仅是不小于。
正文中的线性下界因此是保守的。

\section{命题 \ref{prop:three-grounds} 的三条根据是否穷尽}

这一论证有限度。三条根据的形式如下：

\begin{enumerate}[label=(g\arabic*)]
\item 偏好是生成性的：存在意图 $I$ 与情形 $s$，使得 $I$ 的某个成分只有在意图主体
      遭遇 $s$ 之后才被确定。
\item 规格是有损压缩：存在意图 $I$ 与规格 $\Phi$，使得 $\Phi$ 无法承载 $I$ 的全部
      相关信息。
\item 确认不可归约为验证：存在需求 $R$ 与实现 $P$，使得"$P$ 满足 $\Phi$"为真而
      "$P$ 满足 $R$"为假，且判定后者的过程不能被归约为判定前者。
\end{enumerate}

三条根据不是互相独立的：(g1) 是 (g2) 的一个来源，(g3) 是 (g2) 的一个后果。
因此更准确的表述可能是：\emph{意图鸿沟的根源是 (g1) 与 (g2) 的复合}。本文保留三条
并列的表述，是因为它们在工程上对应不同的缓解手段：(g1) 对应交互式精化与反例生成；
(g2) 对应追溯与覆盖分析；(g3) 对应独立的确认流程与责任划分。

是否存在第四条根据，是一个开放问题（第 \ref{ch:objections} 章 O 项）。

\section{定理 \ref{th:condorcet} 的界是否紧}

式 \eqref{eq:chernoff} 是 Hoeffding 界，它在大 $n$ 下不是紧的。对
$p < 1/2$，多数错误的概率实际上由中心极限定理给出，
$\Pr[\text{多数出错}] \approx \Phi\bigl(-\sqrt{n}\,(1/2-p)/\sigma\bigr)$，
其指数为 $\Theta(n (1/2-p)^{2})$ 但常数更优。

本文选择 Hoeffding 界而非更紧的估计，原因是：

\begin{enumerate}[label=(\alph*)]
\item 它不需要对个体错误的分布作任何假设，只需要独立性、有界性与均值；
\item 它的形式更简单，因而不容易被误读为"精确值"；
\item 本文的核心结论（命题 \ref{th:correlated}）不依赖该界的紧致性：
      相关错误模式给出的下界 $q$ 与上界的常数无关。
\end{enumerate}

\section{命题 \ref{th:correlated} 的一个更强的形式}

正文给出了 $\Pr[\text{多数出错}] \ge q$。可以给出一个更细致的形式。设共同错误
模式 $C$ 与个体独立错误 $I_i$ 满足：验证者 $i$ 出错当且仅当 $C = 1$ 或
$I_i = 1$。则多数出错的事件为
\[
\{C = 1\} \;\cup\; \Bigl\{C = 0,\; \sum_{i=1}^{n} I_i \ge \tfrac{n}{2}\Bigr\},
\]
故
\[
\Pr[\text{多数出错}] \;=\; q + (1-q)\,\Pr\Bigl[\sum_{i=1}^n I_i \ge \tfrac{n}{2}\Bigr]
\;\ge\; q .
\]
并且当 $q > 0$ 固定时，
$\lim_{n \to \infty} \Pr[\text{多数出错}] = q$。这给出了一个比"下界为 $q$"更强的
陈述：\emph{共识的可靠性存在一个由相关错误模式决定的渐近地板，增加验证者数量只能
逼近该地板，不能突破它}。这正是推论 \ref{cor:diversity} 的定量依据。

\section{命题 \ref{prop:active} 的说明}

正文的论证使用了信息增益的定义。这里补充一个说明：命题 \ref{prop:active} 并不是说
被动阅读没有价值。被动阅读提供的是\emph{先验}（在贝叶斯意义上），而主动探测提供的是
\emph{似然}。一个没有先验的探测者不知道去探测什么；一个只有先验的阅读者无法更新。
因此第 \ref{sec:protocol} 节的七步协议前四步是阅读（建立先验），后三步是探测
（获取梯度）。二者的次序不是任意的。
